\RequirePackage{filecontents}
\begin{filecontents}{DPI_refs.bib}
@article{Aspden2016,
  author = {Aspden, Reuben S. and Padgett, Miles J. and Spalding, Gabriel C.},
  title = {Video recording true single-photon double-slit interference},
  journal = {American Journal of Physics},
  year = {2016},
  volume = {84},
  pages = {671--676}
}
@article{Kolenderski2014,
  author = {Kolenderski, Piotr and others},
  title = {Time-resolved double-slit interference pattern measurement with entangled photons},
  journal = {Scientific Reports},
  year = {2014},
  volume = {4},
  pages = {4685}
}
@article{Iskhakov2016,
  author = {Iskhakov, T. Sh. and others},
  title = {Generation and direct imaging of spatial modes in parametric down-conversion},
  journal = {Optics Express},
  year = {2016},
  volume = {24},
  pages = {14336--14347}
}
@article{Spasibko2017,
  author = {Spasibko, K. Yu. and others},
  title = {Multiphoton spatial mode engineering},
  journal = {Nature Photonics},
  year = {2017},
  volume = {11},
  pages = {709--713}
}
@article{Sharapova2020,
  author = {Sharapova, P. R. and others},
  title = {Spatial mode engineering in spontaneous parametric down-conversion},
  journal = {Physical Review Research},
  year = {2020},
  volume = {2},
  pages = {033495}
}
@article{Gleiser1994,
  author = {Gleiser, M.},
  title = {Pseudostable field configurations in the early Universe},
  journal = {Physical Review D},
  year = {1994},
  volume = {49},
  pages = {2978--2981}
}
@article{Amin2012,
  author = {Amin, Mustafa A. and Hertzberg, Mark P. and Kaiser, David I.},
  title = {Nonlinear Dynamics of Oscillons},
  journal = {Physical Review Letters},
  year = {2012},
  volume = {108},
  pages = {241302}
}
@article{Fodor2006,
  author = {Fodor, G. and Forg\'acs, P. and Grandcl\'ement, P. and R\'acz, I.},
  title = {Oscillons in scalar field theories},
  journal = {Physical Review D},
  year = {2006},
  volume = {74},
  pages = {105018}
}
@book{Goodman2005,
  author = {Goodman, J. W.},
  title = {Introduction to Fourier Optics},
  publisher = {Roberts \& Company},
  year = {2005},
  edition = {3rd}
}
\end{filecontents}

\documentclass[aps,pra,reprint,superscriptaddress,10pt]{revtex4-2}

% Packages
\usepackage{amsmath,amssymb}
\usepackage{graphicx}
\usepackage{dcolumn}
\usepackage[colorlinks=true,linkcolor=blue,citecolor=blue,urlcolor=blue]{hyperref}
\usepackage{siunitx}
\usepackage{bm}

% Metadata
\begin{document}

\title{Detector-Plane Imaging (DPI): Camera-Back Physics and Imprint-Field Dynamics}
\author{Srikar R.}
\affiliation{Department / Institute (placeholder)}
\date{December 11, 2025}

\begin{abstract}
Detector-Plane Imaging (DPI) is introduced as an imaging modality that directly visualizes the spatiotemporal coherence patterns of a nonlinear scalar field on the detector plane. Unlike traditional optical measurements that treat the detector as a histogramming device, DPI interprets the sensor as the final physical boundary of the field, recording intensity patterns that encode coherence, decoherence, and structural evolution. Using a $\phi^4$ nonlinear field model we simulate imprint-field dynamics and demonstrate: (1) coherent ring formation and oscillon breathing, (2) photon-like accumulated intensity buildup, (3) coherence-dependent double-slit far-field structure, and (4) quantitative diagnostics including correlation length $\xi(t)$, oscillon width $\sigma(t)$, spectral structure $S(\mathbf{k})$, and far-field cross-sections. We show that modern cameras already capture DPI-like patterns in experiments involving high-gain PDC, OPOs, and multimode nonlinear optics. DPI provides a unified, measurement-plane perspective on coherence evolution and offers an actionable framework for designing next-generation coherent imaging systems.
\end{abstract}

\maketitle

% ============================
% Sections
% ============================

\section{Introduction}
The detector plane is typically treated as the end of an experiment: a surface that collects light, counts photons, and displays accumulated intensity. DPI reframes this surface as an imaging boundary where pre-measurement structure is directly visible. When a nonlinear field evolves and interacts with a detector-plane geometry, the resulting intensity map carries signatures of coherence length, symmetry, interference capability, and temporal stability.

In standard double-slit experiments, only the accumulated interference fringes are shown. The camera's raw frame-by-frame data—ringing, breathing, speckle evolution—rarely enters the scientific interpretation. DPI highlights that these evolving patterns are not noise but structured field behaviour and develops diagnostics to quantify coherence as a dynamical process.

\section{Methods}
\subsection{Field model}
We evolve a two-dimensional $\phi^4$ nonlinear scalar field on a square lattice. The governing continuum equation (in units where $c=1$) is
\begin{equation}
\frac{\partial^2 \phi}{\partial t^2} = \nabla^2 \phi + m^2 \phi - \lambda \phi^3,
\label{eq:phi4}
\end{equation}
with $m^2=-1.0$ and $\lambda\approx 0.59$ chosen to produce long-lived, localized oscillon solutions from Gaussian initial data. The lattice discretization uses a second-order finite difference Laplacian with grid spacing $\Delta x = 0.5$ on a $256\times256$ domain and a leapfrog-like update with time step $\Delta t = 0.05$.

\subsection{Detector-plane intensity and accumulation}
The detector-plane intensity is defined as
\begin{equation}
I(x,y,t) = \phi(x,y,t)^2,
\end{equation}
and accumulated intensity over $N$ frames is
\begin{equation}
I_{\mathrm{acc}}(x,y;N) = \frac{1}{N}\sum_{n=1}^{N} I(x,y,t_n).
\end{equation}
This mimics multi-frame photon accumulation and allows direct comparison to single-photon buildup experiments.

\subsection{Double-slit mask and far-field}
We apply a vertical double-slit transmission mask $M(x,y)$ to the detector-plane intensity and compute the far-field intensity via the Fourier transform:
\begin{equation}
F(\mathbf{k}) = \mathcal{F}\{I(x,y)\cdot M(x,y)\},\qquad I_{\mathrm{ff}}(\mathbf{k})=|F(\mathbf{k})|^2,
\end{equation}
where $\mathcal{F}$ denotes a discrete FFT. Cross-sections of $I_{\mathrm{ff}}$ provide a quantitative measure of fringe visibility.

\subsection{Coherence diagnostics}
We define metrics:
\begin{itemize}
  \item correlation length $\xi(t)$ from the $1/e$ decay of the radial Fourier power around $k=0$;
  \item oscillon width $\sigma(t)=\sqrt{\langle r^2\rangle}$ computed from intensity-weighted second moment;
  \item spectral map $S(\mathbf{k})=|{\cal F}\{I\}|^2$ (log-scaled for display);
  \item normalized far-field cross-sections for visibility quantification.
\end{itemize}

\section{Results}
\subsection{Imprint-field evolution}
DPI simulations reveal four main stages: (i) early coherent rings; (ii) peak coherence with strong ring contrast; (iii) onset of decoherence and speckle emergence; and (iv) slow decay of spatial correlation. These stages map directly onto detector-plane frames and are visible in single-shot and accumulated intensity sequences.

\begin{figure}[tb]
  \centering
  \includegraphics[width=\columnwidth]{Figure1_ImprintEvolution.png}
  \caption{Imprint-field evolution on the detector plane. (a) Early coherent rings, (b) peak coherence, (c) onset of decoherence, (d) pseudo-3D intensity proxy. These frames are sampled from the simulated $\phi^4$ evolution and plotted as $I=\phi^2$.}
  \label{fig:imprint_evolution}
\end{figure}

\subsection{Accumulated intensity buildup}
Accumulated frames show convergence toward a smooth intensity envelope as the number of frames increases, analogous to single-photon interference buildup observed experimentally.

\begin{figure}[tb]
  \centering
  \includegraphics[width=\columnwidth]{Figure2_AccumulatedIntensity.png}
  \caption{Accumulated intensity analogues: progressive accumulation over representative frame counts shows emergence of the stable imprint-field structure.}
  \label{fig:accumulation}
\end{figure}

\subsection{Double-slit illumination and far-field}
Applying a double-slit mask to coherent imprint illumination yields a double-band far-field signature with high fringe visibility. Under decoherent illumination the band structure washes out and visibility decreases.

\begin{figure}[tb]
  \centering
  \includegraphics[width=\columnwidth]{Figure3_DoubleSlit_Farfield.png}
  \caption{Double-slit DPI results. Left: coherent detector-plane illumination. Center: slit aperture. Right: far-field intensity (log scale) showing double-band structure for coherent illumination and degraded structure for decoherent illumination.}
  \label{fig:doubleslit}
\end{figure}

\subsection{Quantitative diagnostics}
We compute $\xi(t)$ and $\sigma(t)$ along with spectral maps $S(\mathbf{k})$ and normalized far-field cross-sections. These diagnostics quantify the onset and progression of decoherence and provide replicable measures for experimental verification.

\begin{figure}[tb]
  \centering
  \includegraphics[width=\columnwidth]{Figure4_Diagnostics.png}
  \caption{Quantitative diagnostics: (a) correlation length $\xi(t)$, (b) oscillon width $\sigma(t)$, (c) Fourier-space intensity $S(\mathbf{k})$ at peak coherence, (d) far-field cross-sections normalized for comparison between coherent and partially decoherent illumination.}
  \label{fig:diagnostics}
\end{figure}

\section{Photon--Oscillon Correspondence}
Oscillons in the $\phi^4$ field act as localized coherent excitations whose imprint-field dynamics resemble the cumulative behavior of photons in interference experiments. This is an operational correspondence: both oscillons and photon fields produce structured detector-plane intensity maps whose evolution encodes coherence and decoherence dynamics. Importantly, DPI does not claim photons are literally oscillons; the mapping is phenomenological and diagnostic.

\section{Experimental Realization}
DPI is immediately testable with extant imaging hardware. We list representative instrument classes that can realize DPI measurements:

\begin{itemize}
  \item Intensified sCMOS / ICCD (Andor iStar, PI-MAX, Hamamatsu ORCA-Quest): single-shot or gated imaging with nanosecond resolution.
  \item EMCCD (Andor iXon series, Princeton ProEM): low-light and single-photon accumulation.
  \item SPAD arrays (SwissSPAD2, MegaFrame): time-tagged single-photon detection with picosecond timing.
  \item High-speed CMOS sensors (Raspberry Pi HQ, modern smartphone sensors): bright-source demonstrations.
\end{itemize}

\noindent\textbf{Minimal DPI test protocol:}
(1) place the camera at the detector plane (no imaging optics); (2) illuminate with a structured coherent source (high-gain PDC, OPO, or engineered laser profile); (3) record short-exposure and accumulation sequences; (4) optionally insert a double-slit mask and compare far-field patterns between coherent and decoherent illumination.

\section{Discussion}
Detector-Plane Imaging (DPI) reframes the role of the detector in optical and field-theoretic experiments. Rather than treating the camera as a passive endpoint that registers discrete photon events, DPI shows that the detector plane itself contains a dynamically accessible record of field coherence. The sequence of ring formation, breathing, fragmentation, and eventual decoherence observed in the imprint field is not a computational artifact: it is a direct diagnostic of internal synchrony within the evolving nonlinear scalar field.

The oscillon analogy demonstrates that coherence-driven spatial ordering emerges naturally in nonlinear dynamics without invoking quantum mechanical assumptions. Yet the same spatial signatures---structured interference, visibility decay, and cumulative fringe buildup---appear in laboratory experiments traditionally interpreted through single-photon language. Heralded-photon double-slit experiments (Aspden \textit{et al.}, Kolenderski \textit{et al}.) show the same gradual progression from granular noise to high-visibility interference that DPI reproduces in purely classical field simulations. This alignment suggests that a significant portion of what we call "wave--particle duality" reflects the detector-plane manifestation of underlying coherence, regardless of whether the source is classical, nonlinear, or quantum.

Importantly, DPI provides a way to quantify decoherence as a continuous dynamical process through correlation-length decay $\xi(t)$, spectral narrowing $S(\mathbf{k})$, and aperture-filtered visibility transitions. These diagnostics exist at the level of the field---not the particle record---and therefore generalize across experimental platforms. Modern cameras (ICCD, EMCCD, SPAD, sCMOS) already record the raw spatiotemporal features DPI identifies; the framework presented here merely provides the missing interpretation.

DPI thus establishes a measurement-plane perspective in which coherence, not particle identity, dictates the morphology of observed interference patterns. This suggests a broader conceptual shift: imaging coherence directly at the detector surface may be as fundamental as counting individual detection events. The method points naturally toward a new class of coherence-aware instruments capable of resolving dynamical pre-measurement structure rather than only final detection outcomes.

\section{Conclusion}
This manuscript introduced Detector-Plane Imaging (DPI), a measurement-plane approach that maps the formation and decay of pre-measurement coherence in a nonlinear scalar field. Through oscillon-driven imprint dynamics, DPI demonstrated that emergent field structures can act as robust coherence sources capable of producing high-contrast interference patterns. The resulting behavior aligns closely with time-resolved single-photon interference experiments, indicating that many features typically attributed to wave--particle duality are already encoded in the coherence structure visible directly at the detector plane.

By interpreting decoherence as a continuous dynamical process---rather than an abrupt measurement-induced collapse---DPI positions the detector as an active coherence imager. This conceptual shift provides a foundation for a more general framework of Quantum Microscopy, in which spatial coherence is treated as a measurable physical field. The results presented here motivate the development of new sensor architectures designed to capture pre-measurement coherence, opening pathways toward enhanced imaging, improved diagnostics, and a deeper understanding of dynamical field behavior in both classical and quantum systems.

\section*{Acknowledgements}
The author thanks the open-source JAX, NumPy, and Matplotlib communities, whose tools enabled the high-resolution simulations used in this work. The manuscript benefited from comparisons to publicly available single-photon datasets and optical imaging studies. No external funding was used for this research, and all computations were performed using publicly accessible hardware and open-source tools.

\bibliographystyle{apsrev4-2}
\bibliography{DPI_refs}

\end{document}
